\chapter{Stationary Action and the Euler--Lagrange Equation}

\lectureoverview{The familiar name \emph{principle of least action} hides the
essential point: the physical trajectory makes the action stationary, but need
not minimize it.  We shall build that statement from ordinary equilibrium,
extend it from points to entire functions, derive the Euler--Lagrange equation
by discretizing a path, and recover Newton's equations.  Rotating and polar
coordinates will then show why the Lagrangian formulation is more than a
change of notation.}

\section{Stationary, Not Necessarily Least}

The word \emph{least} is the unreliable part of the traditional name.  A
classical trajectory does not always make the action a minimum.  It may instead
give a maximum or a saddle in the space of possible trajectories.  The correct
principle is that the action is \emph{stationary}: its first-order change
vanishes.

Begin with the simpler question of equilibrium.  A particle that remains at
rest has no acceleration, so its net force must vanish.  For a force derived
from a one-dimensional potential,
\begin{equation}
  F(x)=-\frac{dV}{dx},
  \label{eq:force-from-v}
\end{equation}
the equilibrium condition is
\begin{equation}
  \frac{dV}{dx}=0.
  \label{eq:stationary-v-one}
\end{equation}

\begin{figure}[htbp]
  \centering
  \includegraphics[width=0.74\textwidth]{lecture_03_figure_02.png}
  \caption{The one-dimensional equilibrium condition: force vanishes where
  the potential is stationary.}
  \lectureframe{00:02:51}
\end{figure}

Now displace the particle by a small amount \(\delta x\).  To first order,
\begin{equation}
  \delta V
  =V(x+\delta x)-V(x)
  =\frac{dV}{dx}\,\delta x+O(\delta x^2).
  \label{eq:first-variation-one}
\end{equation}
At equilibrium the linear term is zero for either sign of \(\delta x\):
\begin{equation}
  \delta V=0
  \qquad\hbox{to first order.}
  \label{eq:delta-v-zero-one}
\end{equation}
This does not say that \(V\) is exactly unchanged after a finite move.  It
says that the leading change is quadratic or higher.  A minimum and a maximum
both satisfy this condition.

\begin{figure}[htbp]
  \centering
  \begin{tikzpicture}[scale=0.95,>=Latex]
    \draw[->] (-0.2,0) -- (3.35,0) node[right] {$x$};
    \draw[->] (0,-0.2) -- (0,2.25) node[above] {$V$};
    \draw[blue!65!black,thick,smooth]
      plot coordinates {(0.2,1.85) (0.75,0.85) (1.55,0.35)
                        (2.35,0.85) (3.0,1.85)};
    \fill[red!65!black] (1.55,0.35) circle (2pt);
    \node at (1.55,-0.34) {minimum};
    \begin{scope}[xshift=4.25cm]
      \draw[->] (-0.2,0) -- (3.35,0) node[right] {$x$};
      \draw[->] (0,-0.2) -- (0,2.25) node[above] {$V$};
      \draw[blue!65!black,thick,smooth]
        plot coordinates {(0.2,0.35) (0.75,1.25) (1.55,1.85)
                          (2.35,1.25) (3.0,0.35)};
      \fill[red!65!black] (1.55,1.85) circle (2pt);
      \node at (1.55,-0.34) {maximum};
    \end{scope}
  \end{tikzpicture}
  \caption{Stationarity fixes the first derivative, not whether the point is
  a minimum or a maximum.}
  \label{fig:min-max-stationary}
\end{figure}

For many coordinates \(q_i\), the first variation is
\begin{equation}
  \delta V
  =\sum_i\frac{\partial V}{\partial q_i}\,\delta q_i.
  \label{eq:first-variation-many}
\end{equation}
If this must vanish for every independent infinitesimal displacement
\(\delta q_i\), then every coefficient vanishes:
\begin{equation}
  \frac{\partial V}{\partial q_i}=0
  \qquad\hbox{for every }i.
  \label{eq:stationary-v-many}
\end{equation}

A two-dimensional potential is better imagined as hilly terrain than as a
roller coaster.  A stationary point may be a valley bottom, a hilltop, or a
saddle.  At a saddle, motion in one direction initially climbs while motion in
another initially descends.  The restoring behavior can therefore be stable
in one direction and unstable in another, even though all first derivatives
vanish at the point itself.

\begin{figure}[htbp]
  \centering
  \includegraphics[width=0.76\textwidth]{lecture_03_figure_06.png}
  \caption{A saddle is stationary in every first-order direction, although
  neighboring displacements rise along one axis and fall along another.}
  \lectureframe{00:08:20}
\end{figure}

In a space with many coordinates, saddle points are more typical than pure
maxima or minima.  This is why \emph{stationary} is the durable word.  The
condition
\begin{equation}
  \delta V=0
  \label{eq:equilibrium-variation}
\end{equation}
summarizes equilibrium without prejudging the second-order behavior.

\section{From Points to Functions}

Mechanics concerns whole motions, not isolated equilibrium points.  A
trajectory is a function such as \(x(t)\), or a collection
\((x(t),y(t),z(t))\).  The mathematical problem of making a quantity
stationary by varying an entire function is the \emph{calculus of variations}.
Before applying it to motion, consider three simpler examples.

\subsection{Shortest curves and geodesics}

Fix two points in the plane and represent a curve between them by \(y(x)\).
Over a short segment,
\begin{equation}
  ds=\sqrt{dx^2+dy^2}
  =dx\sqrt{1+\left(\frac{dy}{dx}\right)^2},
  \label{eq:line-element}
\end{equation}
where the last form assumes \(dx>0\).  The total length is the functional
\begin{equation}
  S[y]=\int_{x_0}^{x_1}
  \sqrt{1+y'(x)^2}\,dx.
  \label{eq:arc-length-functional}
\end{equation}
Square brackets emphasize that \(S\) takes an entire function as input.  The
straight line is selected by
\begin{equation}
  \delta S=0
  \label{eq:stationary-length}
\end{equation}
under small changes of the curve that leave both endpoints fixed.

\begin{figure}[htbp]
  \centering
  \includegraphics[width=0.78\textwidth]{lecture_03_figure_07.png}
  \caption{A trial curve between two fixed endpoints.  Its length is the sum
  of the infinitesimal hypotenuses \(ds\).}
  \lectureframe{00:13:36}
\end{figure}

The same question can be asked on a curved surface.  Pulling a string tight
between two stakes gives a physical picture of the shortest constrained
route.  Such a stationary-length curve is a \emph{geodesic}.  On a flat plane
it is a straight line; on a sphere or hilly surface it generally is not.

\subsection{Fermat's principle}

A light ray in a medium with position-dependent speed \(c(x,y)\) need not
follow the shortest spatial route.  The time required for a small segment is
distance divided by local speed, so
\begin{equation}
  T[y]
  =\int_{x_0}^{x_1}
  \frac{\sqrt{1+y'(x)^2}}{c(x,y)}\,dx.
  \label{eq:fermat-functional}
\end{equation}
Fermat's principle requires
\begin{equation}
  \delta T=0.
  \label{eq:fermat-stationary}
\end{equation}
When \(c\) is constant this reduces to the shortest-distance problem.  In a
variable refractive medium, however, the fastest path balances distance
against local propagation speed.

\begin{figure}[htbp]
  \centering
  \includegraphics[width=0.80\textwidth]{lecture_03_figure_08.png}
  \caption{The least-time functional for a ray in a medium whose local light
  speed depends on position.}
  \lectureframe{00:22:50}
\end{figure}

A topical joke invokes the then-recent claim that neutrinos moved unusually
fast in Italy.  The physics argument immediately returns to the reliable
example: ordinary light changes speed in material according to the index of
refraction.

\subsection{A hanging chain}

A chain pinned at two points settles into a catenary.  Each small piece
contributes gravitational potential energy, so the total potential is itself
an integral over the chain.  Once again the endpoints are fixed, the unknown
shape is a function, and an integral must be stationary.

\begin{figure}[htbp]
  \centering
  \includegraphics[width=0.76\textwidth]{lecture_03_figure_03.png}
  \caption{A hanging chain between fixed supports.  Its stationary shape is
  selected by the total gravitational potential energy.}
  \lectureframe{00:30:41}
\end{figure}

These examples have the same architecture: two fixed endpoints, a function
between them, and an integral whose first variation vanishes.  Reducing such a
global instruction to a differential equation is what it means to solve the
variational problem.

\section{Mechanics as a Boundary-Value Problem}

A trajectory is often drawn with time horizontal because \(t\) is the
independent variable.  There is nevertheless a strong convention in physics
to draw time upward.  An aside recalls experiments in which people
ordered pictures from youth to age in directions associated with their
writing systems: left-to-right readers tended to arrange them left to right,
right-to-left readers the opposite way, and vertically trained readers
vertically.  Physicists, the joke concludes, are the people who insist that
time flows upward.  We shall use the conventional horizontal \(t\)-axis here
because it makes \(x(t)\) read naturally as a graph.

In the initial-value formulation, the data at \(t_0\) are
\begin{equation}
  x(t_0)=x_0,
  \qquad
  \dot x(t_0)=v_0.
  \label{eq:initial-value-data}
\end{equation}
Newton's equation then determines the subsequent orbit.  A variational
formulation poses a different problem: specify two endpoint events,
\begin{equation}
  x(t_0)=x_0,
  \qquad
  x(t_1)=x_1,
  \label{eq:boundary-value-data}
\end{equation}
and determine the physical path between them.  The initial velocity is no
longer supplied; it becomes part of the solution.

\begin{figure}[htbp]
  \centering
  \includegraphics[width=0.76\textwidth]{lecture_03_figure_09.png}
  \caption{A mechanical trajectory \(x(t)\).  One may prescribe position and
  slope at one end, or positions at two endpoint times.}
  \lectureframe{00:27:41}
\end{figure}

The quantity to be made stationary is the \emph{action},
\begin{equation}
  A[x]=\int_{t_0}^{t_1}L(x,\dot x,t)\,dt.
  \label{eq:action-one}
\end{equation}
The integrand \(L\) is the \emph{Lagrangian}.  It may depend on position,
velocity, and explicitly on time.  The action principle is
\begin{equation}
  \delta A=0
  \label{eq:stationary-action}
\end{equation}
for all small path variations with
\begin{equation}
  \delta x(t_0)=\delta x(t_1)=0.
  \label{eq:fixed-endpoint-variations}
\end{equation}
The endpoint condition is essential: we compare neighboring paths that connect
the same two events.

\section{A Discrete Derivation of the Euler--Lagrange Equation}

The action is a functional of infinitely many path values.  To make its
variation concrete, divide the time interval into steps of length
\(\epsilon\),
\begin{equation}
  t_n=t_0+n\epsilon,
  \qquad
  x_n=x(t_n),
  \label{eq:discrete-times}
\end{equation}
and approximate the velocity on the \(n\)-th interval by
\begin{equation}
  v_n=\frac{x_{n+1}-x_n}{\epsilon}.
  \label{eq:discrete-velocity}
\end{equation}
The integral becomes the ordinary sum
\begin{equation}
  A_\epsilon
  =\sum_{n=0}^{N-1}
  \epsilon L\left(x_n,\frac{x_{n+1}-x_n}{\epsilon}\right).
  \label{eq:discrete-action}
\end{equation}
We use \(n\) for the time-slice label so that \(i\) remains available later
for spatial coordinates and particles.

\begin{figure}[htbp]
  \centering
  \includegraphics[width=0.82\textwidth]{lecture_03_figure_10.png}
  \caption{The trajectory replaced by a chain of time slices.  Varying one
  interior height changes only its two neighboring segments.}
  \lectureframe{00:38:30}
\end{figure}

Vary one interior point \(x_n\), holding every other \(x_k\) fixed.  Only two
terms in the sum depend on \(x_n\):
\begin{align}
  A_\epsilon\big|_{x_n}
  &=
  \epsilon L\left(x_n,\frac{x_{n+1}-x_n}{\epsilon}\right)
  \notag\\
  &\quad+
  \epsilon L\left(x_{n-1},
                  \frac{x_n-x_{n-1}}{\epsilon}\right).
  \label{eq:two-neighbor-terms}
\end{align}
Differentiating gives
\begin{align}
  0=\frac{\partial A_\epsilon}{\partial x_n}
  &=
  \epsilon\frac{\partial L_n}{\partial x_n}
  -\frac{\partial L_n}{\partial v_n}
  +\frac{\partial L_{n-1}}{\partial v_{n-1}}.
  \label{eq:discrete-el-variation}
\end{align}
The opposite signs have a simple geometric origin.  Raising \(x_n\) increases
the slope of the interval behind it and decreases the slope of the interval
ahead of it.

\begin{figure}[htbp]
  \centering
  \includegraphics[width=0.74\textwidth]{lecture_03_figure_04.png}
  \caption{The direct position derivative and the two neighboring
  velocity-derivative contributions in the discrete variation.}
  \lectureframe{00:45:26}
\end{figure}

Rearrange Eq.~\eqref{eq:discrete-el-variation}:
\begin{equation}
  \frac{1}{\epsilon}
  \left(
    \frac{\partial L_n}{\partial v_n}
    -\frac{\partial L_{n-1}}{\partial v_{n-1}}
  \right)
  =
  \frac{\partial L_n}{\partial x_n}.
  \label{eq:discrete-el}
\end{equation}
The left side is a neighboring difference quotient.  As
\(\epsilon\to0\), the broken line becomes a smooth path and the quotient
becomes a time derivative:
\begin{equation}
  \boxed{
  \frac{d}{dt}\left(\frac{\partial L}{\partial\dot x}\right)
  =
  \frac{\partial L}{\partial x}}
  .
  \label{eq:euler-lagrange-one}
\end{equation}
This is the Euler--Lagrange equation.

\begin{figure}[htbp]
  \centering
  \includegraphics[width=0.78\textwidth]{lecture_03_figure_11.png}
  \caption{The completed Euler--Lagrange equation obtained from stationarity
  at every interior time slice.}
  \lectureframe{00:49:25}
\end{figure}

Once \(L\) is specified, Eq.~\eqref{eq:euler-lagrange-one} acts as a machine:
differentiate the given function, and the machine returns a differential
equation for the path.

\begin{classroomqa}
  \audiencequestion{Is the Euler--Lagrange result valid for any function of
  \(x\) and \(\dot x\)?\lecturetimestamp{00:49:37--00:50:24}}
  \lecturerresponse{Yes, subject to the regularity needed for the derivatives.
  The arc-length example has exactly the same form after relabeling its
  independent variable \(x\) as the role played by \(t\), and its unknown
  function \(y(x)\) as the role played by \(x(t)\).  Its Euler--Lagrange
  equation has straight-line solutions.}
\end{classroomqa}

\section{Newton's Equations from the Lagrangian}

For a particle on a line in a potential \(V(x)\), take
\begin{equation}
  L(x,\dot x)=T-V
  =\frac12m\dot x^2-V(x).
  \label{eq:lagrangian-tv-one}
\end{equation}
The minus sign may look strange because total energy is \(T+V\).  But the
Lagrangian is not the energy.  Its test is whether the Euler--Lagrange machine
returns the correct motion.  A mock vote about the sign is declined: the
answer must be checked, not elected.

First,
\begin{equation}
  \frac{\partial L}{\partial\dot x}=m\dot x.
  \label{eq:conjugate-momentum-one}
\end{equation}
This is momentum.  More generally, the momentum conjugate to any coordinate
\(q\) is defined by
\begin{equation}
  p_q=\frac{\partial L}{\partial\dot q}.
  \label{eq:conjugate-momentum-general}
\end{equation}
Taking the time derivative of Eq.~\eqref{eq:conjugate-momentum-one} gives
\(m\ddot x\), while
\begin{equation}
  \frac{\partial L}{\partial x}
  =-\frac{dV}{dx}.
  \label{eq:lagrangian-position-derivative}
\end{equation}
The Euler--Lagrange equation is therefore
\begin{equation}
  m\ddot x=-\frac{dV}{dx}=F(x),
  \label{eq:newton-from-action}
\end{equation}
which is Newton's second law.

\begin{figure}[htbp]
  \centering
  \includegraphics[width=0.80\textwidth]{lecture_03_figure_05.png}
  \caption{Kinetic energy, potential energy, \(L=T-V\), and the recovered
  component form of Newton's equation on one board.}
  \lectureframe{00:57:52}
\end{figure}

For many coordinates \(q_i\),
\begin{equation}
  L(\mathbf q,\dot{\mathbf q})
  =\sum_i\frac12m_i\dot q_i^{\,2}-V(q_1,q_2,\ldots),
  \label{eq:lagrangian-tv-many}
\end{equation}
and there is one equation for each coordinate:
\begin{equation}
  \frac{d}{dt}\left(\frac{\partial L}{\partial\dot q_i}\right)
  =\frac{\partial L}{\partial q_i}
  \quad\Longrightarrow\quad
  m_i\ddot q_i=-\frac{\partial V}{\partial q_i}=F_i.
  \label{eq:euler-lagrange-many}
\end{equation}
The trajectory is still one curve, now through a high-dimensional
configuration space.

\begin{classroomqa}
  \audiencequestion{Have we merely worked backward from Newton's law and
  engineered a functional that reproduces the answer?\lecturetimestamp{00:58:07--01:00:04}}
  \lecturerresponse{That is not the productive viewpoint.  The Lagrangian is
  the compact specification of the physical system: one function of all
  coordinates and velocities packages all its equations of motion.  In new
  theories, physicists usually propose an action and derive the equations,
  rather than guessing every equation separately.}
\end{classroomqa}

This packaging is one advantage.  Coordinate freedom is another.  The
statement that a line is shortest does not depend on whether the plane is
described by Cartesian, polar, or curvilinear coordinates.  In the same way,
stationarity of the action is coordinate-independent, even though the formula
for \(L\) and the resulting component equations change.

\section{A Free Particle Seen from a Carousel}

Let \((x,y)\) be inertial Cartesian coordinates and let \((X,Y)\) rotate
uniformly with angular velocity \(\omega\).  With the orientation used here,
\begin{align}
  x&=X\cos\omega t+Y\sin\omega t,
  \notag\\
  y&=-X\sin\omega t+Y\cos\omega t.
  \label{eq:rotating-coordinate-map}
\end{align}
The angle between the frames is \(\omega t\).

\begin{figure}[htbp]
  \centering
  \includegraphics[width=0.76\textwidth]{lecture_03_figure_12.png}
  \caption{Inertial and uniformly rotating Cartesian axes.  The same point
  receives time-dependent rotating coordinates.}
  \lectureframe{01:06:05}
\end{figure}

A free particle has inertial-frame Lagrangian
\begin{equation}
  L_{\mathrm{in}}
  =\frac{m}{2}\left(\dot x^2+\dot y^2\right).
  \label{eq:free-inertial-lagrangian}
\end{equation}
Differentiate Eq.~\eqref{eq:rotating-coordinate-map} and use
\(\sin^2\omega t+\cos^2\omega t=1\).  The trigonometric terms combine to give
\begin{align}
  \dot x^2+\dot y^2
  &=
  \dot X^2+\dot Y^2
  +\omega^2(X^2+Y^2)
  +2\omega(Y\dot X-X\dot Y).
  \label{eq:rotating-speed}
\end{align}
Thus
\begin{align}
  L_{\mathrm{rot}}
  &=
  \frac{m}{2}(\dot X^2+\dot Y^2)
  +\frac{m\omega^2}{2}(X^2+Y^2)
  +m\omega(Y\dot X-X\dot Y).
  \label{eq:rotating-lagrangian}
\end{align}

\begin{figure}[htbp]
  \centering
  \includegraphics[width=0.84\textwidth]{lecture_03_figure_13.png}
  \caption{The transformed free-particle Lagrangian: ordinary kinetic energy,
  a quadratic centrifugal term, and a term linear in velocity.}
  \lectureframe{01:12:16}
\end{figure}

The quadratic coordinate term has the form \(-V_{\mathrm{fict}}\), where
\begin{equation}
  V_{\mathrm{fict}}
  =-\frac{m\omega^2}{2}(X^2+Y^2).
  \label{eq:centrifugal-potential}
\end{equation}
It decreases outward, so it produces the centrifugal force
\(m\omega^2(X,Y)\).

The final term in Eq.~\eqref{eq:rotating-lagrangian} is linear in both
coordinates and velocities.  Applying the Euler--Lagrange equation gives
\begin{align}
  m\ddot X&=m\omega^2X-2m\omega\dot Y,
  \notag\\
  m\ddot Y&=m\omega^2Y+2m\omega\dot X.
  \label{eq:rotating-equations}
\end{align}
The terms proportional to \(\omega^2\) are centrifugal; those proportional to
velocity are Coriolis.  In vector form within the rotating plane,
\begin{equation}
  \mathbf F_{\mathrm{C}}
  =-2m\,\boldsymbol\omega\times\mathbf v_{\mathrm{rot}},
  \label{eq:coriolis-vector}
\end{equation}
with signs adjusted if the opposite orientation convention is chosen.

\begin{figure}[htbp]
  \centering
  \includegraphics[width=0.84\textwidth]{lecture_03_figure_14.png}
  \caption{The rotating-frame equations.  Coordinate-proportional terms give
  centrifugal acceleration; cross-velocity terms give Coriolis acceleration.}
  \lectureframe{01:20:20}
\end{figure}

The point is not that the algebra disappears.  It is that one coordinate
transformation of a single scalar function generates both coupled equations
and both fictitious forces.

\begin{classroomqa}
  \audiencequestion{Does this method work for arbitrarily curved
  coordinates?\lecturetimestamp{01:22:01--01:22:17}}
  \lecturerresponse{Yes.  The new coordinates may be curvilinear and may even
  depend explicitly on time, as the rotating coordinates do.  Rewrite the
  action in the new variables and apply the same Euler--Lagrange rule.}
\end{classroomqa}

\section{Polar Coordinates and a Cyclic Angle}

For ordinary polar coordinates,
\begin{equation}
  x=r\cos\theta,
  \qquad
  y=r\sin\theta.
  \label{eq:polar-map}
\end{equation}
Both \(r\) and \(\theta\) may vary with time.  Differentiating and adding the
squares yields
\begin{equation}
  \dot x^2+\dot y^2
  =\dot r^2+r^2\dot\theta^2.
  \label{eq:polar-speed}
\end{equation}
For a central potential \(V(r)\),
\begin{equation}
  L
  =\frac{m}{2}\left(\dot r^2+r^2\dot\theta^2\right)-V(r).
  \label{eq:polar-lagrangian}
\end{equation}

\begin{figure}[htbp]
  \centering
  \includegraphics[width=0.80\textwidth]{lecture_03_figure_15.png}
  \caption{Polar geometry and the decomposition of kinetic energy into radial
  and angular motion.}
  \lectureframe{01:23:30}
\end{figure}

The radial equation is
\begin{equation}
  m\ddot r
  =mr\dot\theta^2-\frac{dV}{dr}.
  \label{eq:radial-equation}
\end{equation}
The first term is the outward centrifugal contribution associated with the
particle's own angular motion.

For the angle,
\begin{equation}
  \frac{d}{dt}\left(\frac{\partial L}{\partial\dot\theta}\right)
  =\frac{\partial L}{\partial\theta}=0,
  \label{eq:theta-el}
\end{equation}
because the central potential has no explicit \(\theta\)-dependence.  Hence
\begin{equation}
  \frac{d}{dt}\left(mr^2\dot\theta\right)=0,
  \qquad
  mr^2\dot\theta=\ell,
  \label{eq:angular-momentum}
\end{equation}
where \(\ell\) is the conserved angular momentum.

\begin{figure}[htbp]
  \centering
  \includegraphics[width=0.82\textwidth]{lecture_03_figure_16.png}
  \caption{The angular Euler--Lagrange equation has zero right-hand side,
  exposing conservation of \(mr^2\dot\theta\).}
  \lectureframe{01:29:25}
\end{figure}

The angular equation gives
\begin{equation}
  \dot\theta=\frac{\ell}{mr^2}.
  \label{eq:theta-dot-angular-momentum}
\end{equation}
Substitution into Eq.~\eqref{eq:radial-equation} leaves a single equation for
\(r\):
\begin{equation}
  m\ddot r
  =\frac{\ell^2}{mr^3}-\frac{dV}{dr}.
  \label{eq:effective-radial-equation}
\end{equation}

A coordinate absent from \(L\) is traditionally called \emph{cyclic}, even
when it is not an angle.  If \(q\) is cyclic, then
\begin{equation}
  \frac{\partial L}{\partial q}=0
  \quad\Longrightarrow\quad
  \frac{d}{dt}\left(\frac{\partial L}{\partial\dot q}\right)=0.
  \label{eq:cyclic-conservation}
\end{equation}
Its conjugate momentum is conserved.  This is the first glimpse of a general
link between symmetry and conservation law.

\section{Questions About Meaning, Energy, and Scope}

The questions that follow expose where the deeper content of the formalism
lies.

\begin{classroomqa}
  \audiencequestion{Is there an independent reason to arrive at the
  Lagrangian, beyond discovering that it gives the right classical
  equations?\lecturetimestamp{01:33:25--01:35:07}}
  \lecturerresponse{Historically, more than a century of classical mechanics
  showed that particles, waves, and electromagnetic fields could all be cast
  in this form.  At the classical level that accumulated success is an
  empirical fact about how laws are organized.  The deeper explanation comes
  from quantum mechanics and is deferred.}
\end{classroomqa}

\begin{classroomqa}
  \audiencequestion{In the rotating-frame Lagrangian, can the
  velocity-linear term be assigned an ordinary kinetic or potential
  energy?\lecturetimestamp{01:35:08--01:36:23}}
  \lecturerresponse{No.  The conserved quantity associated with time
  translation is obtained by the Hamiltonian construction, not by reading
  every term in \(L\) as kinetic or potential energy.  For
  Eq.~\eqref{eq:rotating-lagrangian}, the velocity-linear term cancels from
  \(\sum_a\dot q_a\,\partial L/\partial\dot q_a-L\).  The Hamiltonian treatment
  will organize this systematically.}
\end{classroomqa}

Indeed, because \(L_{\mathrm{rot}}\) has no explicit time dependence,
\begin{equation}
  H_{\mathrm{rot}}
  =\dot X\frac{\partial L}{\partial\dot X}
   +\dot Y\frac{\partial L}{\partial\dot Y}
   -L
  =\frac{m}{2}(\dot X^2+\dot Y^2)
   -\frac{m\omega^2}{2}(X^2+Y^2)
  \label{eq:rotating-hamiltonian-preview}
\end{equation}
is conserved.  It is the rotating-frame Jacobi integral, not the same
numerical energy assigned by the inertial observer.  Time-dependent coordinate
changes need not preserve the numerical form of energy.

\begin{classroomqa}
  \audiencequestion{What if a curved trajectory is caused by
  gravity?\lecturetimestamp{01:36:30--01:37:10}}
  \lecturerresponse{For Newtonian gravity, simply use
  \(V(r)=-GMm/r\) in the central-force Lagrangian.  Curved spacetime belongs
  to general relativity and is not assumed here.}
\end{classroomqa}

\begin{classroomqa}
  \audiencequestion{Are the centrifugal potential and its energy merely
  fictitious?\lecturetimestamp{01:37:13--01:40:14}}
  \lecturerresponse{An observer on the rotating platform measures the
  effective forces and can consistently use the rotating-frame conserved
  quantity.  \emph{Fictitious} means that these terms arise from a noninertial
  description, not that their observed effects are imaginary.  Energy itself
  is not invariant under a general time-dependent change of coordinates.}
\end{classroomqa}

\begin{classroomqa}
  \audiencequestion{What direction does the velocity-dependent Coriolis
  force point?\lecturetimestamp{01:40:22--01:42:36}}
  \lecturerresponse{It is perpendicular to the rotating-frame velocity, with
  magnitude proportional to speed and to the frame's angular velocity.  In
  the plane the components exchange \(\dot X\) and \(\dot Y\) with opposite
  signs, as in Eq.~\eqref{eq:rotating-equations}.}
\end{classroomqa}

\begin{figure}[htbp]
  \centering
  \includegraphics[width=0.76\textwidth]{lecture_03_figure_17.png}
  \caption{The cross-velocity component pattern: the Coriolis force is
  perpendicular to the rotating-frame velocity.}
  \lectureframe{01:41:40}
\end{figure}

\begin{classroomqa}
  \audiencequestion{Must the starting Lagrangian be written in an inertial
  frame?\lecturetimestamp{01:42:37--01:44:00}}
  \lecturerresponse{For an ordinary free particle, Newton's inertial-frame law
  supplies the simple starting form \(m(\dot x^2+\dot y^2)/2\).  Once that
  physical law is known, the action can be transformed to any other
  coordinates.  One should not assume the same simple kinetic formula
  independently in a rotating frame.}
\end{classroomqa}

The same type of velocity-linear term occurs for a charged particle in a
magnetic field:
\begin{equation}
  L_{\mathrm{em}}
  =\frac12m\mathbf v^2
   +q\mathbf A(\mathbf r,t)\cdot\mathbf v
   -q\phi(\mathbf r,t).
  \label{eq:electromagnetic-lagrangian}
\end{equation}
For a uniform magnetic field perpendicular to the plane, the symmetric gauge
\(\mathbf A=\tfrac12\mathbf B\times\mathbf r\) produces a term proportional to
\(X\dot Y-Y\dot X\), the same algebraic structure as the Coriolis term.

\begin{figure}[htbp]
  \centering
  \includegraphics[width=0.82\textwidth]{lecture_03_figure_18.png}
  \caption{The velocity-linear rotating-frame term that motivates the
  comparison with a charged particle in a uniform magnetic field.}
  \lectureframe{01:44:40}
\end{figure}

\begin{classroomqa}
  \audiencequestion{Are there constraints on what function may serve as a
  Lagrangian, beyond smoothness?\lecturetimestamp{01:44:46--01:45:06}}
  \lecturerresponse{Yes, but the relevant constraints require more structure
  than has yet been developed and are deferred.  At this stage the central
  lesson is how a proposed Lagrangian generates equations and conservation
  laws.}
\end{classroomqa}

The rotating equations also admit circular motion: a point fixed in the
inertial frame is seen from the carousel to circle the rotating origin.  The
magnetic analogy then explains why a charged particle can execute circular
cyclotron motion.

\begin{classroomqa}
  \audiencequestion{Would a charged particle remain on that circle, or would
  radiation make the orbit collapse?\lecturetimestamp{01:46:11--01:47:21}}
  \lecturerresponse{The particle-only equations being discussed omit
  radiation.  Once coupling to the electromagnetic field and radiation
  reaction are included, an accelerating classical charge can radiate.
  Radiation is therefore a further physical effect, not a contradiction of
  the orbit in the truncated equations.}
\end{classroomqa}

\begin{classroomqa}
  \audiencequestion{Where do hurricanes enter this discussion?
  \lecturetimestamp{01:47:30--01:47:50}}
  \lecturerresponse{Their large-scale atmospheric motion occurs on the
  rotating Earth, so Coriolis effects are important.  A proper derivation
  belongs to meteorology and is not attempted here.}
\end{classroomqa}

\begin{classroomqa}
  \audiencequestion{Is the action framework expected to cover every physical
  law, and what would a genuine exception mean?\lecturetimestamp{01:47:57--01:49:01}}
  \lecturerresponse{Every known closed fundamental theory in the discussion's
  scope can be organized by an action: Newtonian particles, Maxwell fields,
  general relativity, and the Standard Model.  Apparent counterexamples such
  as viscous drag arise when environmental degrees of freedom are omitted.
  Enlarge the system to include the fluid or microscopic bath and the
  fundamental description again fits the framework.  A truly fundamental
  exception would therefore be a major discovery.}
\end{classroomqa}
